% WORKING DRAFT v2 of the MARL paper (30/09/2026, from 16:0x EDT).  Replaces the selector framing of paper_marl.tex,
% which is kept unchanged as the historical draft (with its A1/G1 corrections of 30/09).
% Part I (analysis + minimal fix) rests on evidence in hand.  Part II (a theory of prefix-broadcast communication) is IN
% PROGRESS: its mechanism checks are measured, its E3-specific statements are NOT yet in Lean (rule: nothing enters as a
% theorem before compiled, gated Lean).  Every number is copied from a named ledger row of
% research/papers/audits/MARL_SELF_AUDIT_LEDGER_20260928.md (tags L-xxx in the notes) and its artifacts.
% All E3 results are DEVELOPMENT or DEVELOPMENT-CONFIRM (pre-registered, unused panels); the reserved test 1301-1316 is closed.
% Build: tectonic -X compile paper_marl_v2.tex --keep-logs
\documentclass[11pt]{article}
\usepackage[margin=1in]{geometry}
\usepackage{amsmath,amssymb,amsfonts,amsthm}
\usepackage{graphicx,booktabs,url,xcolor}
\usepackage[hidelinks]{hyperref}

\newtheorem{proposition}{Proposition}
\newtheorem{conjecture}{Conjecture}
\theoremstyle{definition}
\newtheorem{definition}{Definition}
\newtheorem{prediction}{Prediction}

\newcommand{\lean}[1]{\textsf{\footnotesize #1}}
\newcommand{\PROVED}{\textsc{proved}}
\newcommand{\NUM}{\textsc{numerical}}
\newcommand{\DEV}{\textsc{dev}}
\newcommand{\PRE}{\textsc{pre-reg}}
\newif\ifshownotes\shownotestrue
\definecolor{notecol}{RGB}{0,70,190}
\definecolor{todocol}{RGB}{190,20,20}
\newcommand{\NOTE}[1]{\ifshownotes{\color{notecol}\,\textsf{\footnotesize[N: #1]}}\fi}
\newcommand{\TODO}[1]{\ifshownotes{\color{todocol}\,\textsf{\footnotesize[TODO: #1]}}\fi}
\newcommand{\TODOPAR}[1]{\ifshownotes\par\vspace{2pt}\noindent%
  {\color{todocol}\textsf{\footnotesize\textbf{To decide.}\;#1}}\par\vspace{2pt}\fi}

\begin{document}
\sloppy
\title{Silence Is a Strong Baseline: Relay, Rotation and the Distilled Zero-Byte Bar\\
for Prefix-Broadcast Multi-Agent Communication}
\author{\TODO{authors and affiliation}}
\date{Working draft v2, 30 September 2026}
\maketitle

\TODOPAR{(1) Venue and page limit are unconfirmed (the v1 draft named CCSC Eastern); they decide how much of
Sections~\ref{sec:fail}--\ref{sec:theory} moves to an appendix. (2) Framing: Part I (analysis plus minimal fix) is
supported now; Part II (Section~\ref{sec:theory}) is in progress and must stay labelled as such until its E3-specific
statements are compiled in Lean. (3) Whether to run the reserved test (traffic seeds 1301--1316) on relay with a rotating
head. (4) More training seeds: the gains depend on the teacher's quality (Section~\ref{sec:limits}).}

\begin{abstract}
When several agents share one broadcast stream and each receiver can decode only a prefix of it, of a size that varies
between receivers and over time, does communication help a learned team? On a $3\times3$ CityFlow traffic-signal grid
with such a channel we find that the answer depends on the comparator. Learned actors that use the delivered messages
are worse than an actor trained without any messages. A student distilled from a full-information teacher that sends
\emph{zero} bytes is better than both. Controls show that the student's gain is the teacher's information, not
distillation or extra compute: a student of a silent teacher gains nothing on three seeds, and the gain grows with the
teacher's information. Measured against this distilled silence, five pre-registered communication designs inside the
packet channel fail, and one has a seed-dependent result. We trace the failures to three mechanisms of prefix decoding:
sparse, varying sender sets in a slot-indexed input; stale cached copies; and a congestion-ranked stream head that never
injects low-ranked senders when capacity is low. The minimal repair keeps the channel and its byte accounting. Every
25-byte packet carries a 3-byte decision summary of its sender and aged relays of every other agent in an explicit
185-bit layout, and the head of the stream rotates. On an unused, pre-registered panel this beats the distilled
zero-byte student on both training seeds under every capacity condition tested ($-2.2$ to $-4.0\%$ total waiting). Each
missing piece breaks it: no relay, no rotation, reinforcement learning instead of distillation, or no messages at all.
A companion analysis (in progress) ties each failure to a theorem shape from compression and experimental design:
frozen-surrogate starvation, one ranking for every budget, quota tracking, decision lifetime and obstruction floors. It
checks the predicted mechanisms directly.
\end{abstract}

\part*{Part I --- Analysis and a minimal fix}

\section{Introduction}\label{sec:intro}
Budgeted multi-agent communication is usually evaluated against agents trained without messages. We study a channel
where that comparator is too weak. One broadcast stream is committed per step; after it is committed, each receiver
learns how many packets it can decode, and it decodes the first $K_i$ of the stream. Capacities differ between receivers
and change every step. The same structure appears in layered multicast~\cite{rlm}, where heterogeneous receivers
subscribe to a prefix of layers. It is new, as far as we can tell, as a multi-agent reinforcement-learning channel.

We report three findings and one repair:
\begin{enumerate}
\item \textbf{Messages harm, and silence is strong.} Actors trained on the channel's own deliveries are worse than a
co-trained actor that ignores communication. The heuristic reference used in earlier drafts (max pressure) is itself
harmed by its neighbour messages, and a message-free max-queue rule is the honest heuristic (Section~\ref{sec:fail}).
\item \textbf{The bar is a distilled zero-byte student.} DAgger distillation of a full-information teacher into a policy
that receives no messages beats the co-trained silent actor on five of five seeds. The gain is information, not
distillation or compute (Section~\ref{sec:student}). Against this bar every design inside the packet channel that we
pre-registered fails or is seed-dependent.
\item \textbf{The failure has mechanisms.} We identify three (Section~\ref{sec:mech}): sparse and varying sender sets in a
slot-indexed input, stale cached copies, and head starvation when capacity is low.
\item \textbf{A minimal repair works.} Relay tables in the unchanged packets, plus a rotating stream head, beat the bar
on an unused pre-registered panel, on both seeds and under every capacity condition (Section~\ref{sec:fix}).
\end{enumerate}

\paragraph{Contributions, scoped.}
(i) The prefix-broadcast channel with heterogeneous receiver capacities as a MARL testbed, with explicit byte accounting
and an oracle-access audit. (ii) The distilled zero-byte student as a required comparator for deployment-communication
claims, with controls for distillation, compute and dose. (iii) A diagnosis of why learned communication fails on this
channel, with each alternative hypothesis tested by a rule written before its result. (iv) Relay with a rotating head as
a minimal, pre-registered repair. What is \emph{not} claimed: relay itself is prior art (Section~\ref{sec:related}), and
so is distillation of a privileged teacher.

\section{Related work}\label{sec:related}
\paragraph{Learned communication under budgets.} Scheduling who talks (SchedNet~\cite{schednet}), sparse or compressed
messages (NDQ~\cite{ndq}, IMAC~\cite{imac}, information-bottleneck and vector-quantised messages~\cite{ibvq}),
event-triggered sending~\cite{etcnet} and request-based perception (Where2comm~\cite{where2comm}) are established.
Chauhan and Bera~\cite{chauhan2026} find that learned protocols at 2-bit budgets are often indistinguishable from
silence and that engineered senders beat learned ones. Our results agree in spirit. We add the distilled-silence
comparator, a channel with prefix decoding, and a positive design.
\paragraph{Multi-hop and delayed information.} Delayed multi-hop information structures, in which agent $i$ holds
$x_j(t-k)$ from $k$-hop neighbours, are learned by imitation of a centralised expert in graph neural controllers (Gama et
al.~\cite{gama2022}; Tolstaya et al.~\cite{tolstaya2019}). NeurComm~\cite{neurcomm} propagates delayed global information
through recurrent messages in traffic-signal control. Message chains (TEM~\cite{tem}), two-hop relays (AC2C~\cite{ac2c})
and small-diameter topologies with memory (ExpoComm~\cite{expocomm}) spread information further. Our relay tables are an
instance of this idea. What differs is the channel: none of these models a prefix-decoded stream with heterogeneous
capacities, and none compares with a distilled silent student.
\paragraph{Privileged teachers and distillation.} Distilling centralised or privileged teachers into decentralised
students is established: CTDS~\cite{ctds}, PTDE~\cite{ptde}, communication-enhanced tutoring~\cite{tutoring},
KD-MARL~\cite{kdmarl}, and MAGPO~\cite{magpo}. We use it as a comparator and as the training route, not as a
contribution.
\paragraph{One order for every budget; freshness.} One ordering that serves every cardinality is studied in incremental
maximisation~\cite{bernstein2022,bienkowski2026}. Embedding priorities so that any prefix decodes the important part is
Priority Encoding Transmission~\cite{pet}. Keeping every source fresh on a shared channel is age-of-information
scheduling~\cite{kadota2018}. One sequence that serves every house size within quota is house-monotone apportionment
(Balinski and Young~\cite{balinski1982}).
\paragraph{Traffic.} CityFlow~\cite{cityflow}; CoLight~\cite{colight}; UniComm~\cite{unicomm} (sharing predictions of
approaching vehicles); max pressure~\cite{varaiya2013}; MAPPO~\cite{mappo}; DAgger~\cite{dagger}.

\section{Setting}\label{sec:setting}
\paragraph{Channel (PROTOCOL\_HETERO\_v0).} The $m=9$ intersections are senders and receivers. At each of 26 decision
steps (10\,s apart) the scheduler commits an order $\pi$ of the senders. Receiver capacities $K_i$ are then revealed,
and receiver $i$ decodes $S_i$, the first $K_i$ senders of $\pi$ other than itself. We use three capacity conditions: iid
$K_i\sim U\{1,\dots,7\}$ per receiver and step; \emph{cap4}, the stream capped at 4 packets; and \emph{low},
$K_i\sim U\{1,2,3\}$. A frontier study uses a constant $K$. Packets are 25\,B, and a profile is charged for its stream
prefix plus its description, exactly as in the protocol. The average step costs 188.7\,B under iid, 100\,B under cap4
and about 94\,B under low. A receiver keeps its last direct copy of every sender.
\paragraph{Information.} The scheduler may use the payloads paid in earlier steps, their measurement times, a mirror of
every receiver's cache, the public topology and the actor code. Fresh states, counterfactual views and oracle rollouts
are used only as training labels or labelled diagnostics. Every evaluation runs a per-step access audit, and every
oracle call inside a deployable path raises.
\paragraph{Task.} CityFlow $3\times3$ grid, headways 8 and 16\,s. The metric is total queue (waiting) per episode.
Development panel: traffic seeds 901--964 $\times$ both headways (128 episodes). Pre-registered panels use unused seed
blocks between 2001 and 2664, each used once. Comparisons are paired by episode, reported as the ratio of summed waiting
minus one, with 95\% bootstrap intervals over traffic seeds (both headways resampled together, 4000 draws).

\section{Actors, training and comparators}\label{sec:methods}
\paragraph{Actors.} A flat actor concatenates eight sender blocks (fresh flag, heard flag, age, phase, elapsed,
vehicles and waiting per approach, static topology) with the receiver's own state into an MLP. A \emph{set} actor applies
one shared encoder to each present sender block and mean-pools over present senders, so absent senders contribute
nothing. Both are trained by MAPPO with a centralised critic~\cite{mappo}.
\paragraph{Teacher and students.} The \emph{teacher} is a set actor trained with full, fresh delivery (every receiver
decodes all 8 senders every step). A \emph{student} is trained by DAgger~\cite{dagger}: it acts on its own deployment
input, the teacher labels its states from the counterfactual full view (a training-only oracle), and the teacher mixture
$\beta$ decays from 1 to 0 over the first half of 500 iterations.
\paragraph{Comparators.}
\begin{itemize}
\item \emph{maxq}: serve the approach with the largest queue; 0 bytes.
\item \emph{set\_none}: a set actor co-trained with no messages.
\item \emph{The zero-byte student}: the teacher distilled into a policy with no message input.
\item \emph{FO}: the teacher deployed with a fresh-only receiver, which drops stale cached blocks.
\item \emph{BYTEBUDGET}, a reference in a different channel model: receiver byte capacity $25K_i$, with 3\,B packets, so
that every receiver decodes all 8 senders at 27\,B per step.
\end{itemize}
\paragraph{Relay packets.} Sender $k$'s 25\,B packet carries positional slots:
\begin{itemize}
\item its own 17-bit summary: phase 4, elapsed 1, vehicles $2\times4$, waiting $1\times4$ bits;
\item for each other agent $j$, the same 17 bits of the freshest summary of $j$ that $k$ holds, plus a 4-bit age in
decision steps, saturated at 14, with 15 meaning absent.
\end{itemize}
That is $17+8\times21=185$ bits $\le 200$. Summary levels come from a codebook fitted on training traffic. A receiver
merges the packets it decodes with its own cache, keeping the freshest copy per agent.
\NOTE{L-162, L-165 (the byte-layout audit), L-174 (the saturated-age re-measure).}
\paragraph{Stream order.}
\begin{itemize}
\item \emph{Prefix}: the scheduler's cached-queue ranking (the default).
\item \emph{Agefirst}: senders whose last paid packet is oldest go first.
\end{itemize}
Both use scheduler information only. The cap rule is applied under every order.\NOTE{L-168: a tooling bug that ignored
cap4 under agefirst was found and fixed before any affected run.}

\section{Failures: messages harm, and why}\label{sec:fail}
\paragraph{Learned actors that use their messages are worse than silence.} On the development panel, a flat actor
trained with the same recipe but no messages has $8.0\%$ $[7.4,8.6]$ less waiting than the committed-prefix system of
five flat prefix-trained actors, and less than each of them ($5.6$ to $10.5\%$). A second training seed gives $7.9\%$.
Doubling the prefix actor's training to 4000 iterations leaves it $6.3\%$ $[5.3,7.4]$ worse than the silent actor at
equal iterations.\NOTE{L-105, L-109, L-113.}
\paragraph{The heuristic reference was harmed by its own messages.} Max pressure as implemented has about 14\% MORE
waiting than the same rule without its downstream (message-fed) terms. This max-queue rule uses no messages, and every
pressure variant tried is worse still. Against max-queue, warm-started actors cloned from max pressure are worse
($+1.2$ to $+5.5\%$), and flat prefix actors are only $1.7$--$6.3\%$ better.\NOTE{L-114--L-116.}
\paragraph{Selecting messages does not fix it.} An earlier line of work scheduled messages for frozen prefix-trained
actors. No legal selector beat the committed prefix in closed loop (Appendix~\ref{app:history}). The oracle headroom
behind that effort is, to first order, best-of-N over action perturbations of an imperfect actor: forcing one agent's
action with no message change reaches the same oracle.\NOTE{L-78, L-81.}

\section{Mechanisms of the failure}\label{sec:mech}
We wrote each hypothesis and its reading rule before its result.
\paragraph{Sparse, varying sender sets in a slot-indexed input.}
\begin{itemize}
\item A flat actor trained with full delivery keeps 79--85\% of its full-information gain when 5--7 of 8 senders are
present.
\item It falls below the silent actor under the protocol's capacity process, where about half the senders are absent:
$+4.2$ and $+1.3\%$ on two seeds.
\item Mean pooling over present senders (the set actor) degrades gracefully; attention pooling equals mean pooling.
\item Three alternatives are refuted by their rules: irrelevant input dimensions (H1), stale content alone (H5), and
weight sharing without pooling (H3); doubling the width (H7) does not close the gap.
\item Under 25\% packet loss the set receiver keeps its advantage, and the flat one breaks.
\end{itemize}
\NOTE{L-132, L-133, L-135, L-142, L-157.}
\paragraph{Stale cached copies.}
\begin{itemize}
\item The full-information set teacher, deployed under the prefix with its stale cached blocks, is no better than
set\_none: $-0.1\%$ $[-1.0,+0.9]$.
\item Dropping stale blocks at the receiver (FO) recovers 59\% of the full-information gain: $-3.3\%$ on the development
panel, $-3.0\%$ and $-5.2\%$ for two seeds on a pre-registered panel.
\end{itemize}
\NOTE{L-124, L-125, L-127, L-138.}
\paragraph{Head starvation at low capacity.}
\begin{itemize}
\item The default order ranks senders by cached congestion. When few packets are decoded, the same senders stay at the
head.
\item In a traffic-free replica of the channel, validated against measured ages, a sticky ranking at cap4 leaves 40--46\%
of (receiver, agent) pairs never covered.
\item With relay packets in CityFlow the pairs are still covered, but copies are old: the mean relayed age is 5.4 decision
steps, and 45\% of copies are at least 4 steps old, under cap4 and low.
\item A rotating head gives ages of 0.8--1.4 steps, with 0--5\% of copies at least 4 steps old.
\end{itemize}
\NOTE{L-168 (simulator), L-177 (probe ages).}

\section{The bar: a distilled zero-byte student}\label{sec:student}
\paragraph{Gain.} Development panel, five training seeds: the zero-byte student has $3.5\%$ $[2.2,4.9]$ less waiting than
set\_none, and less on every seed.\NOTE{L-160; two-level bootstrap over training and traffic seeds.}
\paragraph{Controls.}
\begin{itemize}
\item \emph{Distillation alone:} the same DAgger recipe with a silent teacher gives $+0.2$, $-0.2$ and $+0.2\%$ against
set\_none on three seeds.
\item \emph{Extra compute:} 500 more PPO iterations for set\_none give $+2.0\%$.
\item \emph{Dose:} the gain grows with the teacher's information. Against set\_none, a silent teacher gives $\approx0$, a
prefix-message teacher $-2.3/-0.9\%$, and a full-information teacher $-3.2/-6.3\%$ (two seeds).
\end{itemize}
The gain is therefore communication used at training time and compiled into a local policy.\NOTE{L-147, L-149, L-152,
L-153, L-158.}
\paragraph{Designs inside the packet channel measured against the bar} (pre-registered, unused panels, both seeds
required):
\begin{center}\small
\begin{tabular}{@{}lccl@{}}
\toprule
design & seed 0 & seed 1 & verdict\\
\midrule
FO vs set\_none & $-3.0\%$ & $-5.2\%$ & passes against the weak bar\\
FO vs the zero-byte student & $+1.3\%$ & $+2.1\%$ & the student is better\\
fresh-only student vs the zero-byte student & $-0.8\%$ (fail) & $-4.0\%$ & seed-dependent\\
set\_fresh vs the zero-byte student & $-2.0\%$ & $+7.2\%$ (fail) & seed-dependent\\
\bottomrule
\end{tabular}
\end{center}
Extra seeds confirm the reading: set\_fresh is typically about $-2\%$ against set\_none and does not beat the student.
\NOTE{L-127, L-129, L-138, L-140, L-141, L-144--L-146, L-150.}

\section{The fix: relay tables and a rotating head}\label{sec:fix}
\paragraph{Relay removes sparse presence inside the packet channel.} On the development panel, the relay student under
the default order beats the zero-byte student by $3.3\%$ $[2.5,4.2]$ and $4.0\%$ $[3.1,5.0]$ (iid), with the channel and
its accounting unchanged. That equals the byte-budget reference, a tie. With no messages it is $4.8\%$ worse than the
student, so it uses them.\NOTE{L-163; unchanged under saturated ages, L-174.}
\paragraph{Rotation removes starvation.} Under the default order, relay loses its advantage at low capacity. Seed 0 is
$+0.9\%$ (cap4) and $+0.4\%$ (low) against the student on the development panel. A rotating head restores it:
\begin{itemize}
\item cap4: $-2.3$ / $-4.1\%$ (seed 0 / seed 1);
\item low: $-2.8$ / $-3.0\%$;
\item the iid cost of rotation is small and not significant on these panels: $+0.3$ / $+0.9\%$.
\end{itemize}
\NOTE{L-169, L-174.}
\paragraph{Pre-registered confirmation.} We registered the protocol before any run, on panel 2601--2664 (checked unused).
The primary endpoint is relay with a rotating head against the same-seed zero-byte student, and all six cells (two seeds
$\times$ three capacity conditions) had to have an upper 95\% bound below zero.
\begin{center}\small
\begin{tabular}{@{}lccc@{}}
\toprule
& iid (189\,B/step) & cap4 (100\,B/step) & low ($\approx94$\,B/step)\\
\midrule
seed 0 & $-2.19\%$ $[-3.10,-1.25]$ & $-2.96\%$ $[-3.87,-2.04]$ & $-2.19\%$ $[-2.95,-1.43]$\\
seed 1 & $-4.00\%$ $[-4.71,-3.30]$ & $-3.41\%$ $[-4.09,-2.75]$ & $-2.43\%$ $[-3.20,-1.67]$\\
\midrule
\multicolumn{4}{@{}l}{\emph{secondary, same panel}}\\
relay, default order, seed 0 & $-2.71\%$ & $+0.68\%$ $[-0.12,+1.50]$ & $+0.96\%$ $[+0.10,+1.88]$\\
relay, default order, seed 1 & $-4.29\%$ & $-1.76\%$ & $-1.74\%$\\
\bottomrule
\end{tabular}
\end{center}
The primary endpoint is supported in all six cells. The same panel reproduces the starvation effect of the default order
for seed 0 at low capacity.

On the same panel, the relay system also beats:
\begin{itemize}
\item set\_none by 5.7--10.1\%;
\item max-queue by 22.5--26.2\%.
\end{itemize}
The byte-budget reference, in its own channel, gives $-4.0$ / $-4.5\%$.\NOTE{L-173 (protocol,
PROTOCOL\_E3\_RELAY\_ROTATION\_20260930, sha256 c7467697\dots), L-175, report e3\_relayrot\_panel\_report\_20260930.}
\paragraph{Every missing piece breaks it} (development panel, same-seed zero-byte student as the reference):
\begin{itemize}
\item \emph{No relay} (FO): the student is better (Section~\ref{sec:student}).
\item \emph{No rotation}: loses at low capacity (the rows above).
\item \emph{RL instead of distillation}: MAPPO from scratch on relay inputs is $+2.0$ / $+4.5\%$ worse than the student.
\item \emph{No messages}: $+4.8\%$.
\item \emph{A different rotation rule}: a head chosen for the freshest table (a ``bass'' rule) equals age-first rotation,
so we do not claim it.
\end{itemize}
\NOTE{L-166, L-169.}
\paragraph{More seeds (development).} We distilled students trained under rotation (relay6Cage) for seeds 2 and 3. With
a rotating head, against their zero-byte students (iid / cap4 / low):
\begin{itemize}
\item seed 2: $-3.9$ / $-4.6$ / $-4.1\%$;
\item seed 3, whose teacher is weaker: $-1.0$ / $-1.0$ / $+0.0\%$.
\end{itemize}
Training under rotation does not beat rotating only at evaluation (seed 0 / seed 1: $-2.9$ / $-2.7$ / $-1.0$ and $-3.3$ /
$-3.1$ / $-1.9\%$). \TODO{seeds 4--9 relaunched after the disk incident; L-177.}
\paragraph{Reference: a byte-limited channel.} In the alternative channel model where receiver capacity is counted in
bytes, 3\,B decision summaries from the raw full-information teacher beat the zero-byte student on 5/5 seeds, by $3.6\%$
$[2.7,4.5]$. The effect holds:
\begin{itemize}
\item on three pre-registered panels;
\item under unseen demand (headways 6, 12 and 24\,s);
\item with a codebook refitted on other traffic;
\item and up to 25\% packet loss.
\end{itemize}
With eight seeds the mean is $2.5\%$ $[1.4,3.6]$, with the size tracking the teacher's quality. Training directly on
compressed payloads is fragile: one of three seeds fails.\NOTE{L-131, L-144, L-151, L-156, L-157, L-159, L-160, L-170;
the 10-seed synthesis reported by a parallel session is not yet in the ledger.}

\section{Limitations}\label{sec:limits}
\begin{itemize}
\item \emph{Seeds and teacher quality.} The pre-registered relay result rests on two training seeds. Two development
seeds more give one strong pass and one weak result. Across ten teachers the size of every gain tracks the teacher's own
full-information gain; at one seed the teacher is worse than silence.
\item \emph{Rotation at evaluation.} The pre-registered students were trained under the default order. Rotation costs
0.3--1.0\% at iid, and training under rotation did not remove that cost.
\item \emph{Development-confirm only.} The reserved test (traffic seeds 1301--1316) is closed. Pre-registered panels
are unused before their single use, but the designs were selected on the development panel.
\item \emph{One grid, one simulator,} two headways in training, a fixed 26-step horizon, and one codebook family.
\item \emph{Prior art.} Relay, distillation and compression are known. Our claims concern this channel and this
comparator.
\item \emph{Diagnostics use an oracle.} The counterfactual views in Section~\ref{sec:theory} are labelled diagnostics;
no deployable path calls them.
\end{itemize}

\part*{Part II --- Toward a theory of prefix-broadcast communication (in progress)}

\section{Five mechanisms and their direct checks}\label{sec:theory}
\TODOPAR{Every statement in this section is a \emph{shape} taken from a proved result in another model (the source is
named) and a prediction measured in E3. None is a theorem about E3 yet. The E3-specific Lean modules listed at the end
must be compiled and gated before any of it is stated as a theorem.}

\paragraph{M1. Frozen-surrogate starvation and linear repair} (source: the frozen-surrogate separation of Paper A, proved
in its own model). There, every optimum of a frozen surrogate gives zero mass to the decision-relevant cells, and a
$\beta$-mixture with a certificate design buys information linearly in $\beta$, with no threshold. \emph{E3 reading:} the
cached-queue ranking is a frozen surrogate of ``who matters'', and rotation is the repair.
\begin{itemize}
\item \emph{Information level, traffic-free:} coverage of starved pairs has no threshold in $\beta$. At low capacity the
never-covered share falls from 0.41 to 0.27 at $\beta=0.02$ and to 0 at $\beta=0.25$. The mean age, however, first rises
(present but stale copies) before falling.
\item \emph{Closed loop:} the repair grows with $\beta$ and so does the iid sacrifice. Seed 0 shows no threshold at cap4
($-0.96\%$ $[-1.53,-0.43]$ already at $\beta=0.05$). Seed 1 is flat until $\beta\ge0.5$.
\item \emph{Reading:} coverage without bounded age pays little. The operative quantity is the age bound (M3).
\end{itemize}
\NOTE{L-176, L-177.}

\paragraph{M2. One schedule for every budget: signed versus monotone regimes} (source: the exact menu law under signed
cancellation, STRUCTURAL). With non-negative item values one order degrades gracefully across budgets. With signed values
(extra items can hurt) no single order serves every budget. \emph{Prediction:} the waiting-versus-$K$ frontier is
monotone for receivers that drop or age stale inputs and non-monotone for a stale-kept receiver.
\begin{itemize}
\item \emph{Monotone:} relay with the default order, and FO, on both seeds.
\item \emph{Signed:} the stale-kept set actor is non-monotone on seed 0 (waiting rises from $K=2$ to $K=3$), flat on seed
1, and never beats the zero-byte student.
\item \emph{Exception:} relay with rotation is monotone on seed 1 and has one significant violation on seed 0 ($K=5\to6$),
which is unexplained.
\item Relay with rotation beats the student from $K=2$ (seed 0) and $K=1$ (seed 1); relay with the default order needs
$K\ge3$--4.
\end{itemize}
\NOTE{L-177.}

\paragraph{M3. Quota tracking and bounded age} (source: ADAPTIVE; and Paper C's band, whose closest analogue it names as
house-monotone apportionment). A rotating head is an apportionment of head positions over time. For uniform targets the
largest-deficit rule is round-robin and stays within 1 of quota.
\begin{itemize}
\item \emph{ADAPTIVE's warning, reproduced on the E3 scheduler:} with non-uniform targets, the online largest-deficit rule
exceeds unit quota on 209 of 2000 random targets (worst $3/2$; ADAPTIVE's witness $26/23$ reproduced exactly), while
the Balinski--Young quota method never does.
\item \emph{A new pathology:} rotation with deterministic tie-breaks \emph{phase-locks} with raw one-sender packets at
constant $K=2$, leaving 25\% of pairs never covered. Relay removes it.
\item \emph{Measured in CityFlow:} rotation bounds the relayed age (0--5\% of copies are 4 or more steps old, against
45\% under the default order at low capacity).
\end{itemize}
\NOTE{L-176; ADAPTIVE's trace caption claims a unique maximum at every step, but ties occur at $t=0,4,7$; all 18
tie-break paths still give $26/23$.}

\paragraph{M4. Decision lifetime} (source: the functional certificate lifetime, LIFE2/LIFE3; D6's sharp $2\delta$ margin;
the argmax-stability lemma of the MARL certificate module). If the score gap exceeds $2(\delta+h\nu)$, an $h$-step-old
input yields the same decision. \emph{Direct check} (probe, 12 runs, parity 128/128 with the evaluator):
\begin{itemize}
\item Agreement of the relay-view action with the fresh-complete-information action rises with the fresh margin, from
$\approx0.6$ in the lowest to $\approx1.0$ in the highest quintile.
\item It falls with the relayed age: for example 0.94, 0.87, 0.83, 0.81, 0.80 and 0.78 at ages 1--6 under low capacity.
\item A lifetime certificate fitted on half the traffic seeds is sound on the other half ($\ge0.9997$) but covers only
0--23\% of decisions: valid, and conservative.
\end{itemize}
\NOTE{L-177.}

\paragraph{M5. Obstruction floors: when silence is enough} (source: UNIFIED's obstruction criterion and randomised error
floor $1/|H|$; Paper B's safe-order law; SUPPORT's whole-fibre condition). A zero-byte rule is safe exactly when no
obstruction hides inside a local-observation class. Messages are worth bytes only there.
\begin{itemize}
\item \emph{Direct check:} the zero-byte student disagrees with the teacher on 14--17\% of receiver-steps. There, relay
with rotation agrees with the teacher on 53--70\%, and elsewhere on 92--94\%. Relay's teacher-agreement advantage is
therefore concentrated where local information does not settle the decision.
\item Under the default order at low capacity the advantage vanishes or reverses. This links age (M3, M4) to agreement to
waiting.
\item \emph{Decide, not identify:} 3\,B compression flips the teacher's decision on 1.3--2.3\% of steps, almost all in the
smallest-margin quintile.
\end{itemize}
\NOTE{L-177.}

\paragraph{What must be proved for Part II to become a theorem-level contribution} \TODO{Lean, paused by the authors on
30/09}:
\begin{itemize}
\item[(T-a)] Under prefix decoding with $K_i\ge1$, relay merging and an age-first head with a stated tie rule, every
relayed age is bounded (a finite bound in $m$), and a phase-lock counterexample exists for raw packets.
\item[(T-b)] The lifetime composite: a margin above $2(\delta+H\nu)$, with $H$ the bound of (T-a), implies the relay action
equals the fresh-complete action. This instantiates LIFE2/LIFE3 and the MARL argmax lemma.
\item[(T-c)] Monotonicity of the value of a decoded prefix for fresh-only or aged-relay receivers, under a stated
non-negativity condition, against a signed counterexample for stale-kept receivers.
\end{itemize}
Each needs a satisfiability witness.

\section{Reproducibility}
Every number is produced by a script from saved artifacts whose manifests record source hashes and self-checks. The
self-checks are:
\begin{itemize}
\item a parity check of each new evaluator against the previous one;
\item the scheduler-side age mirror, checked against the merged views at every step;
\item probe parity with the evaluator rows;
\item the simulator's reproduction of ADAPTIVE's witness.
\end{itemize}
Main scripts:
\begin{itemize}
\item \texttt{MARL\_e3\_actor\_distill\_v4/v6.py} (students);
\item \texttt{MARL\_e3\_delivery\_eval\_relay\_satrot.py} and \texttt{\_v4.py} (evaluation);
\item \texttt{models/marl\_relay\_v2.py} (the 185-bit layout);
\item \texttt{models/marl\_delivery\_music.py} and \texttt{\_v2.py} (orders);
\item \texttt{MARL\_e3\_theory\_sim.py} and \texttt{MARL\_e3\_theory\_probe.py} (Part II);
\item the reports \texttt{MARL\_e3\_relayrot\_report.py}, \texttt{MARL\_e3\_theory\_cl\_report.py} and
\texttt{MARL\_e3\_theory\_probe\_report.py}.
\end{itemize}
Protocols: \texttt{research/marl\_budget/PROTOCOL\_E3\_*.md}.

\appendix
\section{History: scheduling messages for frozen actors}\label{app:history}
The first design scheduled which packets each receiver decodes, for actors trained with the committed prefix. It used a
paid one-byte receiver bid and a learned critic over candidate profiles; the v1 draft kept its theory (revision
geometry) and its results. On a fixed max-pressure actor the bid improved held-phase selection on fresh traffic, with a
small closed-loop benefit. On learned actors:
\begin{itemize}
\item no legal selector variant beat the committed prefix in closed loop;
\item nor did guarded, decision-focused or densely labelled variants;
\item nor did a SchedNet-style learned schedule.
\end{itemize}
Two later findings explain why this line could not succeed on this testbed. The frozen actors were themselves beaten by
a silent actor (Section~\ref{sec:fail}). And the oracle headroom that motivated selection was best-of-N action
perturbation, not communication value.\NOTE{paper\_marl.tex (v1, kept), L-69, L-78--L-81.}

\begin{thebibliography}{99}
\bibitem{rlm} S.~McCanne, V.~Jacobson, M.~Vetterli, ``Receiver-driven layered multicast,'' SIGCOMM 1996.
\bibitem{pet} A.~Albanese, J.~Bl\"omer, J.~Edmonds, M.~Luby, M.~Sudan, ``Priority encoding transmission,'' FOCS 1994.
\bibitem{kadota2018} I.~Kadota, A.~Sinha, E.~Uysal-Biyikoglu, R.~Singh, E.~Modiano, ``Scheduling policies for minimizing
age of information in broadcast wireless networks,'' IEEE/ACM Trans.\ Networking, 2018, arXiv:1801.01803.
\bibitem{balinski1982} M.~L.~Balinski, H.~P.~Young, \emph{Fair Representation}, Yale University Press, 1982.
\bibitem{schednet} D.~Kim et al., ``Learning to schedule communication in multi-agent reinforcement learning,'' ICLR 2019.
\bibitem{ndq} T.~Wang, J.~Wang, C.~Zheng, C.~Zhang, ``Learning nearly decomposable value functions via communication
minimization,'' ICLR 2020.
\bibitem{imac} R.~Wang, X.~He, R.~Yu, W.~Qiu, B.~An, Z.~Rabinovich, ``Learning efficient multi-agent communication: an
information bottleneck approach,'' ICML 2020, arXiv:1911.06992.
\bibitem{ibvq} ``Bandwidth-efficient multi-agent communication through information bottleneck and vector quantization,''
arXiv:2602.02035, 2026.
\bibitem{etcnet} ``Event-triggered multi-agent reinforcement learning with communication under limited-bandwidth
constraint,'' arXiv:2010.04978.
\bibitem{where2comm} Y.~Hu et al., ``Where2comm: communication-efficient collaborative perception via spatial confidence
maps,'' NeurIPS 2022.
\bibitem{chauhan2026} M.~Chauhan, A.~Bera, ``Emergence, not bandwidth: physical coupling and the limits of learned
multi-agent communication,'' arXiv:2609.34373, 2026.
\bibitem{gama2022} F.~Gama, Q.~Li, E.~Tolstaya, A.~Prorok, A.~Ribeiro, ``Synthesizing decentralized controllers with graph
neural networks and imitation learning,'' IEEE Trans.\ Signal Processing, 2022, arXiv:2012.14906.
\bibitem{tolstaya2019} E.~Tolstaya, F.~Gama, J.~Paulos, G.~Pappas, V.~Kumar, A.~Ribeiro, ``Learning decentralized
controllers for robot swarms with graph neural networks,'' CoRL 2019.
\bibitem{neurcomm} T.~Chu, S.~Chinchali, S.~Katti, ``Multi-agent reinforcement learning for networked system control,''
ICLR 2020, arXiv:2004.01339.
\bibitem{tem} ``Scalable communication for multi-agent reinforcement learning via transformer-based email mechanism,''
IJCAI 2023, arXiv:2301.01919.
\bibitem{ac2c} X.~Wang, X.~Li, J.~Shao, J.~Zhang, ``AC2C: adaptively controlled two-hop communication for multi-agent
reinforcement learning,'' AAMAS 2023, arXiv:2302.12515.
\bibitem{expocomm} X.~Li et al., ``Exponential topology-enabled scalable communication in multi-agent reinforcement
learning,'' ICLR 2025, arXiv:2502.19717.
\bibitem{ctds} ``CTDS: centralized teacher with decentralized student for multi-agent reinforcement learning,''
arXiv:2203.08412.
\bibitem{ptde} ``PTDE: personalized training with distilled execution for multi-agent reinforcement learning,'' IJCAI
2024, arXiv:2210.08872.
\bibitem{tutoring} ``Communication-enhanced tutoring for efficient decentralized multi-agent reinforcement learning,''
arXiv:2508.13661.
\bibitem{kdmarl} ``KD-MARL: resource-aware knowledge distillation in multi-agent reinforcement learning,''
arXiv:2604.06691.
\bibitem{magpo} ``Multi-agent guided policy optimization,'' ICLR 2026, arXiv:2507.18059.
\bibitem{bernstein2022} A.~Bernstein, Y.~Disser, M.~Gro{\ss}, S.~Himburg, ``General bounds for incremental
maximization,'' Mathematical Programming, 2022.
\bibitem{bienkowski2026} Bienkowski et al., ``Incremental submodular maximization: better than greedy,''
arXiv:2606.28558, 2026.
\bibitem{cityflow} H.~Zhang et al., ``CityFlow: a multi-agent reinforcement learning environment for large scale city
traffic scenario,'' WWW 2019.
\bibitem{colight} H.~Wei et al., ``CoLight: learning network-level cooperation for traffic signal control,'' CIKM 2019.
\bibitem{unicomm} ``Multi-agent reinforcement learning for traffic signal control through universal communication
method,'' IJCAI 2022, arXiv:2204.12190.
\bibitem{varaiya2013} P.~Varaiya, ``Max pressure control of a network of signalized intersections,'' Transportation
Research Part C, 2013.
\bibitem{mappo} C.~Yu et al., ``The surprising effectiveness of PPO in cooperative multi-agent games,'' NeurIPS 2022.
\bibitem{dagger} S.~Ross, G.~Gordon, J.~A.~Bagnell, ``A reduction of imitation learning and structured prediction to
no-regret online learning,'' AISTATS 2011.
\end{thebibliography}
\TODO{verify every bibliography entry against its primary source (authors and venue) before submission; entries without
authors were checked only by title and arXiv id.}
\end{document}
